% ===============================================================
\section{Limitations and Future Work}\label{sec:limitations}
% ===============================================================

Several limitations qualify the results above and define the agenda for future
work.

\noindent\textbf{Synthetic, metadata-derived benchmark.} Route labels are
derived from hop-count and cross-document structure rather than annotated by
legal experts, and the benchmark---including the ambiguity and conversation
stress tests---is machine-generated. Although paraphrasing mitigated
template-based label leakage on the strict split, the distribution remains
synthetic and should be validated against a human-authored set before strong
external-validity claims are made.

\noindent\textbf{Weakly supervised clarify class.} The strict split contains
no gold \texttt{clarify} label, so the Stage~1 clarify class is trained on
only $60$ synthetic templates and is deliberately treated as a proposal that
Stage~2 must confirm. End-to-end clarify routing therefore rests on the
verifier and the ambiguity detector rather than on strongly supervised
classification.

\noindent\textbf{In-conversation routing of self-contained queries.} On the
conversation stress test, clearly answerable control queries embedded in a
dialogue are routed correctly only $10\%$ (dense) and $40\%$ (graph/hybrid)
of the time: conversational context inflates the ambiguity and trigger
signals of queries that need no context, and the resolver's
\texttt{not\_needed} status does not yet suppress those signals. Conversation
route accuracy ($0.725$) is bounded by this misrouting rather than by
clarify detection, and the clarify false-positive rate rises from $1.5\%$ on
the ambiguity-free strict split to $11.5\%$ on the diagnostic's answerable
controls as the ambiguity base rate of the surrounding traffic increases.

\noindent\textbf{Sparse legal-effect relations and explicit-only hybrids.}
Cross-document legal-effect relations (\textsc{AMENDS}, \textsc{REPLACES})
remain comparatively sparse in the graph, and the benchmark covers only
\emph{explicit} hybrid questions that name the documents to combine, not
\emph{implicit} ones that require discovering the link first.

\noindent\textbf{Residual retrieval-identifier mismatch.} Evaluation-time
identifier normalisation to the canonical
\texttt{law::document::article} form is implemented, and unresolvable
references are now flagged rather than fuzzily matched. However, the hybrid
merge path and the router-mediated pipeline still \emph{log} source
identifiers that the normaliser cannot resolve, which drives Pure Hybrid
Hit@$1$ to $0.003$ and forces the Hit@$1$ entries of every router-mediated
configuration to be read as unresolved rather than as true coverage. Until
canonical identifiers are emitted at the retrieval stage and the end-to-end
experiments re-run, retrieval-coverage claims for those configurations remain
provisional.

\noindent\textbf{Aggregate parity, not superiority, over dense retrieval.} On
the dense-dominated benchmark the full router ($0.640$) is statistically
indistinguishable from the Pure Vector baseline ($0.663$) on aggregate
token-F1, with overlapping bootstrap intervals. The router's contribution is
therefore argued on exact match ($+71\%$ over Pure Vector), on the causal
Stage~2 ablation, and on clarification safety---not on mean token-F1. A
reader should not infer that the router improves average lexical answer
quality over a strong dense baseline on this distribution.

\noindent\textbf{Metric validity.} Token-F1 against a gold context rewards
lexical overlap, and the citation-constrained generator quotes legal text
verbatim; the combination produces the class-wise inversion of
Table~\ref{tab:per_class_f1}, in which each pure mechanism scores highest on
the \emph{other} mechanism's query class. Cross-system token-F1 differences
therefore partly measure quotation overlap rather than answer correctness,
and per-class token-F1 cannot adjudicate mechanism appropriateness. The
semantic BERT-F1 column specified in Section~\ref{subsec:metrics} is required
to complete this analysis and remains to be computed over the saved
predictions.

\noindent\textbf{Scope of the chain-of-thought claim.} The
\textit{Oracle~+~Stage~2} ablation isolates a causal $+0.074$~F1 contribution
of Stage~2 reasoning under fixed gold routing, but the estimate derives from
a single generator (\texttt{Qwen3-32B-AWQ}), a single prompt family, and
deterministic decoding on one benchmark; its magnitude may not transfer
across models, prompts, or legal domains.

\noindent\textbf{Statistical reporting gaps.} Headline F1 comparisons carry
bootstrap $95\%$ confidence intervals ($B=1000$), but no paired significance
tests are reported, mean reciprocal rank could not be computed because the
per-sample logs recorded only Hit@$1$/Hit@$3$, and the p50/p95 latencies of
the re-run configurations await extraction from the per-query logs.

\noindent\textbf{Routing-accuracy evaluation under template leakage.} Because
paraphrase augmentation preceded splitting, $33.2\%$ of highly similar query
pairs cross the cross-validation fold boundary, inflating the absolute CV
scores of all learned routers to $\approx 0.99$. We therefore headline the
paraphrase-hardened strict-split accuracy ($0.898$) instead, and leave
template-grouped fold construction as an evaluation improvement for future
work.

\noindent\textbf{No empirical comparison to complexity-based adaptive
retrieval.} Adaptive-RAG~\cite{jeong2024adaptive} is the closest prior method
but is not compared empirically, because it routes among reasoning depths for
general-domain English QA rather than among evidence mechanisms for Vietnamese
legal QA. A faithful re-implementation of its complexity router as an
additional baseline---or an explicit statement of why the settings are not
directly comparable---would strengthen the positioning.

\noindent\textbf{Efficiency and hardware dependence.} A Stage~2 verifier call
is roughly $119\times$ the cost of the full Stage~1 routing path in the
routing-stage microbenchmark, and all latencies depend on local vLLM
hardware; end-to-end latency is dominated by generation rather than by
routing, so the efficiency argument is about negligible routing overhead
rather than end-to-end speed.

Future directions include: a three-level cascade
XGBoost~$\rightarrow$~PhoBERT~$\rightarrow$~LLM to preserve routing accuracy
under the linguistic diversity of naturally phrased queries; emitting
canonical \texttt{law::document::article} identifiers at the retrieval stage
with cross-encoder re-ranking to complete the coverage evaluation;
integration of manually reviewed \texttt{clarify} samples and implicit hybrid
questions into the training set; context-gated suppression of conversational
features whenever the history resolver marks a query as self-contained;
stronger legal-relation extraction for
\textsc{AMENDS} and \textsc{REPLACES}; a complexity-router baseline in the
spirit of Adaptive-RAG; semantic and human evaluation of legal answer
correctness and citation faithfulness; and template-grouped cross-validation
for leakage-free router comparison.